\documentclass[11pt]{article}
\usepackage[T1]{fontenc}
\usepackage[utf8]{inputenc}
\usepackage{lmodern}
\usepackage[margin=1in]{geometry}
\usepackage{amsmath,amssymb,amsthm,mathtools}
\usepackage{microtype}
\usepackage{enumitem}
\usepackage{booktabs,array}
\usepackage[hidelinks]{hyperref}
\usepackage{bookmark}
\hypersetup{pdftitle={A resolution-free proof of the Duistermaat--van der Kallen nonvanishing theorem},pdfauthor={Christopher D. Long}}
\setlength{\emergencystretch}{2em}
\numberwithin{equation}{section}
\newtheorem{theorem}{Theorem}[section]
\newtheorem{lemma}[theorem]{Lemma}
\newtheorem{proposition}[theorem]{Proposition}
\newtheorem{corollary}[theorem]{Corollary}
\theoremstyle{definition}
\newtheorem{definition}[theorem]{Definition}
\newtheorem{standard}[theorem]{Standard input}
\theoremstyle{remark}
\newtheorem{remark}[theorem]{Remark}
\newcommand{\C}{\mathbb C}
\newcommand{\R}{\mathbb R}
\newcommand{\Z}{\mathbb Z}
\newcommand{\N}{\mathbb N}
\newcommand{\T}{\mathbb T}
\newcommand{\Pj}{\mathbb P}
\newcommand{\CT}{\operatorname{CT}}
\newcommand{\Newt}{\operatorname{Newt}}
\newcommand{\supp}{\operatorname{supp}}
\newcommand{\conv}{\operatorname{conv}}
\newcommand{\Span}{\operatorname{span}_{\R}}
\newcommand{\Int}{\operatorname{int}}
\newcommand{\relint}{\operatorname{relint}}
\newcommand{\vol}{\operatorname{vol}}
\newcommand{\mass}{\operatorname{mass}}
\newcommand{\Reg}{\operatorname{Reg}}
\newcommand{\Crit}{\operatorname{Crit}}
\newcommand{\dd}{\mathrm d}
\newcommand{\ii}{\mathrm i}
\newcommand{\torus}{(\C^\times)^d}
\newcommand{\OmegaD}{\Omega_d}
\newcommand{\omegaD}{\omega_d}
\newcommand{\abs}[1]{\left|#1\right|}
\newcommand{\norm}[1]{\left\|#1\right\|}
\newcommand{\pair}[2]{\langle #1,#2\rangle}
\newcommand{\arxiv}[1]{\href{https://arxiv.org/abs/#1}{arXiv:#1}}
\title{A resolution-free proof of the\newline Duistermaat--van der Kallen\newline nonvanishing theorem}
\author{Christopher D. Long\thanks{Email: \href{mailto:galizur@gmail.com}{galizur@gmail.com}.}}
\date{September 17, 2026}
\begin{document}
\maketitle
\begin{center}
\small\textit{Research draft for adversarial mathematical audit. Not formally verified.}
\end{center}
\begin{abstract}
We give a resolution-free semialgebraic proof of the Duistermaat--van der
Kallen theorem: if the Newton polytope of a nonzero complex Laurent
polynomial contains the origin, infinitely many positive powers have
nonzero constant term. A uniform middle-dimensional sublevel estimate
for a restricted holomorphic volume form yields logarithmic-form
integrability by dyadic decomposition. For global transport, we use the
closed embedding $z\mapsto(z,z^{-1})$ and its induced Hermitian metric.
Uniform semialgebraic connecting-path bounds imply finiteness of scalar
asymptotic critical values; an explicit normalized-gradient lift then
has complete transport by a linear-growth estimate. The local residue
construction of Duistermaat and van der Kallen identifies the constant-term
generating function with a fiber period. A separate semialgebraic endpoint
sweep gives a primitive with a power-law bound, contradicting the simple
pole forced by vanishing of all positive constant terms. Infinite
nonvanishing follows by applying minimal nonvanishing to powers.
The standard semialgebraic, integration, and ODE inputs are identified;
an independent Whitney--Thom transport proof is retained in an appendix.
\end{abstract}

\section{Introduction and statements}\label{sec:introduction}
Let $d\ge0$ and
\[
 R_d=\C[x_1^{\pm1},\ldots,x_d^{\pm1}].
\]
For $0\ne f=\sum_{\alpha\in\Z^d}a_\alpha x^\alpha\in R_d$, define
\[
 A=\supp(f)=\{\alpha:a_\alpha\ne0\},\qquad
 \Newt(f)=\conv(A),\qquad
 c_n(f)=\CT(f^n)\quad(n\ge0).
\]
In particular $c_0(f)=1$. For $d=0$, the ring is $\C$ and the exponent
space is the zero-dimensional vector space.

\begin{theorem}[Duistermaat--van der Kallen, nonvanishing]\label{thm:main}
Let $d\ge0$ and $0\ne f\in R_d$. If $0\in\Newt(f)$, then for every
integer $M\ge0$ there is an integer $n>M$ such that $c_n(f)\ne0$.
\end{theorem}

\begin{theorem}[Minimal nonvanishing]\label{thm:minimal}
Let $d\ge0$ and $0\ne f\in R_d$. If $0\in\Newt(f)$, there exists
$n\ge1$ such that $c_n(f)\ne0$.
\end{theorem}

Theorem~\ref{thm:main} follows algebraically from
Theorem~\ref{thm:minimal}, by applying it to $f^{M+1}$; see
Section~\ref{sec:corollaries}. Thus only the minimal statement needs an
analytic proof. Strict separation subsequently gives
\[
 \bigl(\CT(f^n)=0\text{ for every }n\ge1\bigr)
 \ \Longrightarrow\
 \CT(hf^n)=0\text{ for all sufficiently large }n
\]
for every fixed Laurent polynomial $h$. The corresponding assertion for
finite Fourier sums on a compact torus follows by integrating characters.

The original proof is due to Duistermaat and van der Kallen
\cite{DvK}. Kiselman's Question~9.1, recording a question of Passare,
asks for a simpler proof of infinite nonvanishing
\cite{Kiselman}. The argument here is organized around a direct estimate
for a restricted complex differential form. Its local residue construction
and primitive-differentiation identity come from
\cite[Sections 4--5]{DvK}; the new argument to be examined is the
sublevel estimate and its use to control chains that may escape to infinity.
We make no claim of bibliographic priority for individual auxiliary lemmas.

\subsection*{Proof architecture and scope}
After the elementary face reduction, assume that $d\ge1$ and
$0\in\Int\Newt(f)$. Put
\[
 \OmegaD=\dd z_1\wedge\cdots\wedge\dd z_d,\qquad
 \omegaD=\frac{\dd x_1}{x_1}\wedge\cdots\wedge\frac{\dd x_d}{x_d}.
\]
The argument has the following implications:
\begin{align*}
 &\text{uniform finite-multiplicity projection estimate}\\
 &\quad\Longrightarrow\quad
 \int_{E\cap\{|P|\le\varepsilon\}}|\OmegaD|
       \le C\varepsilon^{1/D}\\
 &\quad\Longrightarrow\quad
 \int_{S\cap\{|f|\le R\}}|\omegaD|\le C_{S,f}R^{1/D}\\
 &\quad\Longrightarrow\quad
 \left|\int_a^{2a}J(s)\,\dd s\right|\longrightarrow0
       \quad(a\downarrow0).
\end{align*}
If all positive constant terms vanish, the residue identity forces
$J(s)=-(2\pi\ii)^{d-1}/s$. Its integral over $[a,2a]$ is a fixed nonzero
constant, giving a contradiction.

Global transport is established separately by the implications
\begin{align*}
 &\text{Hardt triviality, compact triangulation, and uniform arc lengths}\\
 &\quad\Longrightarrow\quad
       \text{finitely many scalar asymptotic critical values}\\
 &\quad\Longrightarrow\quad
       \text{a normalized-gradient lift with controlled growth}\\
 &\quad\Longrightarrow\quad
       \text{complete fiber transport outside a finite set}.
\end{align*}
The radius and differential norm here come from the closed embedding
$z\mapsto(z,z^{-1})$, so approaching a deleted coordinate hyperplane counts
as escape to infinity. The resulting smooth flow transports homology
classes; it is not asserted to be semialgebraic. Hardt triviality is used
both in the connecting-path lemma and, later, to choose the semialgebraic
sweep on the final real interval.

The term \emph{resolution-free} is literal: no resolution of singularities
is used. The main proof retains semialgebraic decomposition, Hardt
triviality, compact triangulation, integration and Stokes, together with
ordinary finite-dimensional ODE existence, dependence, and continuation.
It does not invoke Whitney stratification, Thom isotopy, or a generalized
Bertini--Sard theorem. Appendix~\ref{app:whitney-thom} gives an independent
Whitney--Thom proof of the finite-exceptional-set and topological
transport conclusion needed for the period argument. The presentation
works relative to the explicitly stated standard inputs; it is not a
reconstruction of their foundations. The stronger critical-value and
exponential-growth statement in \cite[Theorem 5]{DvK} is not a theorem
of this draft.

\section{Elementary Newton-polytope reductions}\label{sec:newton}
\begin{lemma}[Minimal-face reduction]\label{lem:face}
Suppose $0\in\Newt(f)$. There is a Laurent polynomial $\widetilde f$ in
$r\le d$ variables such that
\[
 \CT(\widetilde f^{\,n})=\CT(f^n)\quad(n\ge0),
\]
and either $r=0$ and $\widetilde f$ is a nonzero constant, or $r\ge1$ and
$0\in\Int\Newt(\widetilde f)$ in $\R^r$.
\end{lemma}
\begin{proof}
Let $Q$ be the smallest face of $\Newt(f)$ containing $0$. Then
$0\in\relint Q$. If $Q$ is proper, choose an exposing linear functional
$\ell$ with $\ell=0$ on $Q$ and $\ell>0$ on $\Newt(f)\setminus Q$.
This is possible because the minimum occurs at $0\in Q$. If $Q$ is the
whole polytope, use $\ell=0$.

Let $f_Q$ consist of the terms of $f$ with exponent in $Q$. A sum of
$n$ support exponents can equal $0$ only if every summand belongs to $Q$:
apply $\ell$ and use nonnegativity. Comparing the full coefficient
expansions, including all coefficients and cancellations, gives
$\CT(f^n)=\CT(f_Q^n)$.

The affine span $L$ of $Q$ is linear because $0\in Q$. Since $Q$ is a
lattice polytope, $L$ is rational and $\Lambda=L\cap\Z^d$ is a lattice
of rank $r=\dim L$. Choose a $\Z$-basis of $\Lambda$ and express $f_Q$
in the corresponding group algebra $\C[\Lambda]\simeq R_r$. The zero
exponent and all products are preserved. If $r=0$, its sole coefficient
is nonzero. Otherwise the image of $Q$ has $0$ in its full-dimensional
interior.
\end{proof}

\begin{lemma}[Powers of Newton polytopes]\label{lem:powers}
For $0\ne f\in R_d$ and $q\ge1$,
\[
 \Newt(f^q)=q\Newt(f).
\]
\end{lemma}
\begin{proof}
Every exponent of $f^q$ is a sum of $q$ exponents of $f$, giving one
inclusion. Let $v$ be a vertex of $\Newt(f)$, and choose a linear
functional uniquely maximized at $v$ on the support. If
$\alpha_1+\cdots+\alpha_q=qv$, equality at this maximum forces
$\alpha_1=\cdots=\alpha_q=v$. Thus the coefficient at $qv$ is
$a_v^q\ne0$. All vertices of $q\Newt(f)$ occur in $f^q$, proving the
opposite inclusion. The singleton case is immediate.
\end{proof}

\begin{lemma}[A unimodular chart at a vertex]\label{lem:chart}
Suppose $d\ge1$ and $0\in\Int\Newt(f)$. After a unimodular monomial
change of variables, which preserves every constant term of every power,
one can write
\begin{equation}\label{eq:vertex-chart}
 f(y)=y_1^{-m_1}\cdots y_d^{-m_d}u(y),
 \qquad m_i\in\Z_{>0},\quad u\in\C[y_1,\ldots,y_d],\quad u(0)\ne0.
\end{equation}
\end{lemma}
\begin{proof}
Choose a vertex $v$ of $\Newt(f)$ and let $C$ be the open cone of real
linear functionals whose minimum on the support occurs uniquely at $v$.
Choose a primitive lattice vector $a_1\in C$, and extend it to a lattice
basis $a_1,b_2,\ldots,b_d$ of the dual exponent lattice. Such a primitive
vector exists by rational density in the nonempty open cone, followed by
division by the gcd of its coordinates.

For sufficiently large integers $N_i$, the vectors
$a_i=b_i+N_i a_1$ lie in $C$: divide by $N_i$ and use openness. These
vectors together with $a_1$ still form a lattice basis. In the resulting
exponent coordinates every support exponent $\alpha$ satisfies
$\pair{a_i}{\alpha}\ge\pair{a_i}{v}$, with simultaneous equality only
at $v$. Since $0$ is an interior point, every nonzero linear functional
takes negative values on the support; hence $\pair{a_i}{v}<0$.
Set $m_i=-\pair{a_i}{v}$. Factoring out $y^{-m}$ gives
\eqref{eq:vertex-chart} with constant coefficient $a_v$.

The exponent change is a lattice automorphism, so it preserves zero
exponents, products, and constant terms. From this point one may regard
$f$ as the transformed Laurent polynomial and use the standard ordered
form in the $y$-coordinates; no unrecorded orientation sign is then
introduced.
\end{proof}

\section{Semialgebraic integration and uniform projections}\label{sec:sa}
Semialgebraic means semialgebraic over $\R$, with arbitrary real
parameters allowed. Complex coordinates are identified with their real
and imaginary parts. Dimension in this section is real dimension.

\subsection{The standard inputs}
The following precise forms of standard results will be used.

\begin{standard}[Semialgebraic decomposition and triviality]\label{std:sa}
Coordinate projections of semialgebraic sets are semialgebraic. A
semialgebraic subset of the real line is a finite union of points and
intervals; in particular, if it is unbounded above it contains a tail.
A semialgebraic set and a finite collection of semialgebraic subsets admit
compatible finite decompositions into Nash manifolds. Semialgebraic
families have uniform bounds on the number of connected components of
their fibers. For a continuous semialgebraic map $g:X\to Y$, there is a
finite semialgebraic partition $Y=\coprod_jY_j$ and semialgebraic
homeomorphisms
\[
 g^{-1}(Y_j)\simeq Y_j\times F_j
\]
commuting with projection to $Y_j$. Compatible triangulations may be
chosen for compact semialgebraic sets. Compact singular homology classes
on semialgebraic manifolds admit finite semialgebraic cycle representatives.
\end{standard}

These are the decomposition, triangulation, and Hardt triviality results
in \cite{BCR,Hardt}; the local-triviality statement is
\cite[Theorem 9.3.2]{BCR}. A noncompact set need not have a finite
triangulation by \emph{closed} simplices. Only compact model fibers, compact cycle representatives,
and compact truncated chains will require that assertion here.

\begin{standard}[Integration and Stokes]\label{std:stokes}
Integration of a smooth form over a finite semialgebraic chain is
compatible with finite subdivision and semialgebraic reparametrization.
A bounded semialgebraic set of dimension $k$ has finite $k$-dimensional
Hausdorff measure. For a compact semialgebraic chain $C$ and a smooth
$(k-1)$-form $\eta$ on a neighborhood of its support,
\[
 \int_C\dd\eta=\int_{\partial C}\eta.
\]
On a smooth semialgebraic manifold, integration of a closed form over
semialgebraic cycles agrees with the usual homology pairing.
\end{standard}

The integration convention and change of variables are described in
\cite[Sections 1--2]{HKT}; Stokes is \cite[Theorem 2.9]{HKT}, extended to
finite chains by subdivision and cancellation of common faces. Compatible
$C^1$ triangulations give another foundation \cite{OS}. The finite-volume
assertion can also be obtained by summing the coordinate-projection
bounds proved below. We use ordinary change of variables on smooth
pieces, not a slicing theorem for currents.

\subsection{The absolute density of a complex form}
\begin{definition}\label{def:density}
Let $S\subset\R^N$ be semialgebraic of dimension at most $k$, and let
$\xi$ be a smooth complex-valued $k$-form on a neighborhood of $S$.
On a smooth $k$-dimensional stratum $M$, let $\tau_M(x)$ be either unit
simple $k$-vector spanning $T_xM$. Define
\begin{equation}\label{eq:absolute-density}
 \int_S|\xi|
 =\sum_{\dim M=k}\int_M|\xi_x(\tau_M(x))|\,\dd\mathcal H^k(x).
\end{equation}
Lower-dimensional strata contribute zero. For $k=0$ this is counting
measure. For a finite chain $C=\sum_jn_jC_j$, the positive quantity used
for estimates is $\sum_j|n_j|\int_{C_j}|\xi|$; we write it as
$\int_{|C|}|\xi|$. The oriented integral is linear in the chain.
\end{definition}

The definition is independent of the chosen smooth decomposition, by
common refinement and the change-of-variables formula. It satisfies
\begin{equation}\label{eq:density-identities}
 |g\xi|=|g|\,|\xi|,\qquad
 |\xi+\zeta|\le|\xi|+|\zeta|,\qquad
 \left|\int_C\xi\right|\le\int_{|C|}|\xi|.
\end{equation}
The first identity concerns the modulus of the \emph{full} restricted
complex form. Expanding into real coordinate wedges is permissible for
upper bounds, not as an alternative definition.

For example, on $E=\{(z,-z):1\le|z|\le2\}$ one has
$(\dd z_1\wedge\dd z_2)|_E=0$, although the sum of the absolute values
of its separate real coordinate terms is nonzero. This cancellation will
be essential in Section~\ref{sec:sublevel}.

\subsection{The uniform projection estimate}
\begin{lemma}[Uniform finite-multiplicity integration]\label{lem:projection}
Let $\Theta$ be semialgebraic and let
$\mathcal E\subset\Theta\times\R^N$ be a semialgebraic family with
$\dim E_\theta\le k$. Suppose
$h_\theta:E_\theta\to\R^k$ is a semialgebraic family of continuous
maps, differentiable on the smooth pieces after decomposition. There is
$N_0<\infty$, depending on this entire family, such that
\begin{equation}\label{eq:projection}
 \int_{E_\theta}|\dd h_{\theta,1}\wedge\cdots\wedge\dd h_{\theta,k}|
 \le N_0\vol_k(h_\theta(E_\theta))
 \quad(\theta\in\Theta).
\end{equation}
An unbounded image is allowed; the right side may be infinite.
\end{lemma}
\begin{proof}
Retain the parameter in the map
\[
 H:\mathcal E\longrightarrow\Theta\times\R^k,
 \qquad H(\theta,x)=(\theta,h_\theta(x)).
\]
By Standard input~\ref{std:sa}, the numbers of connected components of
all its fibers are bounded by one integer $N_0$. In particular the
cardinality of every zero-dimensional fiber is at most $N_0$.

Fix $\theta$. Decompose $E_\theta$ into smooth pieces on which the rank
of $h_\theta$ is constant. Pieces of dimension less than $k$, and pieces
on which the rank is less than $k$, have zero pullback $k$-density. On
rank-$k$ pieces of dimension $k$, the map is locally a diffeomorphism.
The set of target points whose whole fiber has positive dimension has
dimension at most $k-1$; its rank-$k$ inverse images also have measure
zero, by local change of variables. Away from this exceptional set there
are at most $N_0$ preimages in total. Apply ordinary change of variables
on injective smooth pieces, or equivalently the area formula with this
counting multiplicity, and sum to obtain \eqref{eq:projection}.
This is the parameter-uniform version of \cite[Proposition 2.5]{HKT};
the parameter-retaining device also appears in the proof of its
Proposition 2.8.
\end{proof}

\begin{remark}\label{rem:uniform-parameters}
In applications, $\theta$ includes all varying set parameters, polynomial
coefficients, and cutoffs. Parameters need not range in a compact set.
Uniformity follows from one fixed semialgebraic description, not from
compactness of parameter space. Adding finitely many prescribed
semialgebraic inequalities produces another fixed family, to which the
same lemma applies with a possibly different constant.
\end{remark}

\begin{corollary}\label{cor:bounded-volume}
For a semialgebraic family $E_\theta$ of dimension at most $d$ contained
in a fixed bounded box in $\C^d$,
\[
 \sup_\theta\int_{E_\theta}|\OmegaD|<\infty.
\]
More generally, the Euclidean $d$-volumes of a bounded semialgebraic family
of subsets of $\R^N$ of dimension at most $d$ are uniformly bounded.
\end{corollary}
\begin{proof}
Expand $\OmegaD$ into its $2^d$ real coordinate wedge terms and use
\eqref{eq:density-identities} and Lemma~\ref{lem:projection} for the
corresponding coordinate projections. Their images lie in fixed boxes.
For the Euclidean-volume assertion, the coordinate components of a unit
simple $d$-vector have squares summing to $1$, so their absolute values
have sum at least $1$. Sum over all coordinate $d$-projections.
\end{proof}


\subsection{Uniform lengths of connecting paths}
\begin{lemma}[Uniform connecting paths]\label{lem:connecting-paths}
Let $\mathcal A\subset\Theta\times\R^m$ be a fixed semialgebraic family
whose fibers $A_\theta$ are compact and contained in the closed unit ball.
There are integers $N\ge1$ and a constant $L\ge1$, depending only on the
family, such that each fiber has at most $N$ connected components and any
two points in the same component can be joined in that component by a
rectifiable semialgebraic path of length at most $L$.
\end{lemma}
\begin{proof}
The component bound is part of Standard input~\ref{std:sa}. Apply Hardt
triviality to the parameter projection. After a finite partition of
$\Theta$, each nonempty family has a semialgebraic trivialization with
one fixed compact model fiber. Triangulate that model by a finite
simplicial complex $K$, and compose its triangulation with the
trivialization to obtain semialgebraic homeomorphisms
$H_\theta:|K|\to A_\theta$, jointly semialgebraic in $\theta$ and the
model point.

For every simplex $\sigma$ of $K$, fix a vertex $v_\sigma$. Given two
points $p,q$ of the same component of $|K|$, choose simplices containing
them. Join $p$ to the chosen vertex of its simplex by a straight segment,
join the two chosen vertices through the one-skeleton of their connected
component, and join the second vertex to $q$. A connected finite complex
has a connected one-skeleton on its vertices. There are finitely many
pairs of vertices, so these edge paths may be fixed in advance with a
uniform bound on the number of segments. The two endpoint segments add
at most two more.

Under $H_\theta$, each segment is constant or an injectively parametrized
semialgebraic arc in the unit ball. As $\theta$ and the endpoint in a
fixed simplex vary, its image belongs to one semialgebraic family of
sets of dimension at most one. Indeed it is the image of
$\{(\theta,p,t):p\in\sigma,\ 0\le t\le1\}$ under
$(\theta,p,t)\mapsto(\theta,p,H_\theta((1-t)p+tv_\sigma))$.
The fixed edge types give finitely many more such families.
Corollary~\ref{cor:bounded-volume}, in real dimension one, bounds the
one-dimensional Hausdorff measure of all these images uniformly.
The length of an injective arc is the measure of its image. Summing
these bounds over the bounded number of segments, and then taking the
maximum over the finitely many Hardt pieces, gives $L$.

The transported parametrizations need not have bounded derivatives.
Rectifiability is what the argument establishes and uses: each arc can
be parametrized by arclength, and decomposed into smooth pieces for the
ordinary chain rule. No Lipschitz bound for the Hardt homeomorphisms is
assumed.
\end{proof}

\begin{remark}\label{rem:connecting-paths-background}
A more general bounded-family path-length theorem is stated as
\cite[Theorem 3.3]{KOS}, attributed there to Teissier. The proof above
supplies the compact-fiber case needed here directly from the
semialgebraic tools and projection bound already used in this paper.
It does not invoke a resolution-based proof of that theorem.
\end{remark}

\section{A middle-dimensional polynomial sublevel estimate}\label{sec:sublevel}
Fix $d\ge1$, $L\ge1$, a semialgebraic family $E_\theta$ of real dimension
at most $d$ in the box $[-L,L]^{2d}\subset\C^d$, and $D\ge1$.
For a polynomial $P$, write
\[
 \|P\|_{\mathrm{coef}}=\max_\alpha|[z^\alpha]P|.
\]

\begin{theorem}[Uniform sublevel estimate]\label{thm:sublevel}
There is $C=C(\mathcal E,d,D,L)$ such that every holomorphic polynomial
$P\in\C[z_1,\ldots,z_d]$ with $\deg P\le D$ and
$\|P\|_{\mathrm{coef}}=1$ satisfies
\begin{equation}\label{eq:sublevel}
 \int_{E_\theta\cap\{|P|\le\varepsilon\}}|\OmegaD|
 \le C\varepsilon^{1/D}
 \qquad(\theta\in\Theta,\ \varepsilon>0).
\end{equation}
The same estimate, with value zero on the left, holds for
$\varepsilon=0$.
\end{theorem}

\begin{lemma}[The one-derivative estimate]\label{lem:derivative}
For each $i\in\{1,\ldots,d\}$, let
\[
 \Xi_i(P)=\dd P\wedge\dd z_1\wedge\cdots
                   \widehat{\dd z_i}\cdots\wedge\dd z_d.
\]
There is $C_1$, uniform in $\theta,P,\varepsilon,r$ and $i$, such that
\begin{equation}\label{eq:one-derivative}
 \int_{E_\theta\cap\{|P|\le\varepsilon,\ |\partial_iP|\ge r\}}
        |\Xi_i(P)|\le C_1\varepsilon
 \qquad(r>0,\ \varepsilon>0).
\end{equation}
The assertion also holds when the derivative cutoff is omitted.
\end{lemma}
\begin{proof}
Expand $\Xi_i(P)$ into real wedges. Each term is the pullback of real
coordinate volume under a map whose first coordinate is $\Re P$ or
$\Im P$, and whose other coordinates are one real coordinate from each
$z_j$ with $j\ne i$. The image of the indicated set lies in
\[
 [-\varepsilon,\varepsilon]\times[-L,L]^{d-1}.
\]
Apply Lemma~\ref{lem:projection} to the family retaining $\theta$, the
real and imaginary coefficients of $P$, $\varepsilon$, and $r$.
There are only finitely many choices of $i$ and real wedge term.
Their image volumes are at most $2\varepsilon(2L)^{d-1}$. Summing proves
the assertion. The triangle inequality is applied only after evaluating
each complex form on the tangent vector.
\end{proof}

\begin{proof}[Proof of Theorem~\ref{thm:sublevel}]
We induct on the degree bound $D$. Corollary~\ref{cor:bounded-volume}
gives a bound $V_0$ for the integral over all of $E_\theta$. It suffices
to handle $0<\varepsilon<1/2$; enlarging the final constant handles
$\varepsilon\ge1/2$.

Let $H=\max\{1,\sqrt2L\}$ and
\[
 B_D=\sum_{0<|\alpha|_1\le D}H^{|\alpha|_1},\qquad
 c=(2B_D)^{-1}.
\]
If every nonconstant coefficient of $P$ has modulus less than $c$, its
constant coefficient has modulus $1$ and the remaining terms have sum
of moduli at most $1/2$ throughout the box. Thus $|P|\ge1/2$ there, and
the sublevel set under consideration is empty.

Otherwise choose a nonconstant coefficient $[z^\beta]P$ of modulus at
least $c$, and an $i$ with $\beta_i>0$. Put $Q=\partial_iP$ and
$M=\|Q\|_{\mathrm{coef}}$. Distinct monomials remain distinct under
this partial differentiation whenever they survive; hence
\begin{equation}\label{eq:derivative-normalization}
 M\ge\beta_i|[z^\beta]P|\ge c.
\end{equation}
Holomorphy gives the exact identity
\begin{equation}\label{eq:wedge-derivative}
 \Xi_i(P)=(-1)^{i-1}Q\OmegaD.
\end{equation}
By \eqref{eq:density-identities}, on $|Q|\ge r$,
\[
 |\OmegaD|\le r^{-1}|\Xi_i(P)|.
\]
Consequently Lemma~\ref{lem:derivative} bounds the large-derivative part
by $C_1\varepsilon/r$.

If $D=1$, then $Q$ is a constant of modulus at least $c$. Applying the
same argument with $r=c$ gives the bound $C_1\varepsilon/c$.
This establishes the base case, including polynomials of degree zero
through the preceding near-constant case.

Suppose $D\ge2$. The induction hypothesis applies to the normalized
polynomial $Q/M$ of degree at most $D-1$, even when it is constant.
Discarding the condition $|P|\le\varepsilon$ on the small-derivative part,
we obtain
\begin{align}
 \int_{E_\theta\cap\{|P|\le\varepsilon\}}|\OmegaD|
 &\le C_1\frac{\varepsilon}{r}
    +C_{D-1}\left(\frac rM\right)^{1/(D-1)}\notag\\
 &\le C_1\frac{\varepsilon}{r}
    +C_{D-1}c^{-1/(D-1)}r^{1/(D-1)}.
 \label{eq:sublevel-split}
\end{align}
Take $r=\varepsilon^{(D-1)/D}$. Both terms are constant multiples of
$\varepsilon^{1/D}$, proving the induction. All constants depend only
on the fixed family and degree bounds. Finally, let
$\varepsilon\downarrow0$ and use monotonicity of the positive integral
to see that the integral over $P=0$ is zero.
\end{proof}

\begin{remark}
This is a sublevel estimate for a restricted complex top form, not for
Euclidean volume of $E_\theta$. The distinction permits the estimate
when a whole component of $E_\theta$ lies in $P=0$: on that component
the relevant complex form must have zero integral density. No regularity
or transversality of the hypersurface $P=0$ is assumed.
\end{remark}

\section{Logarithmic integrability on the torus}\label{sec:log-integrability}
\begin{theorem}[Middle-dimensional logarithmic integrability]\label{thm:log-integrability}
Let $d\ge1$, let $f\in R_d$ satisfy $0\in\Int\Newt(f)$, and let
$S\subset\torus$ be semialgebraic of real dimension at most $d$.
Choose $b\in\Z_{\ge0}^d$ with $b+\supp(f)\subset\Z_{\ge0}^d$ and put
\[
 D=\max_{\alpha\in\supp(f)}|b+\alpha|_1.
\]
Then $D\ge1$, and there is a constant $C=C(S,f,b)$ such that
\begin{equation}\label{eq:log-integrability}
 \int_{S\cap\{|f|\le R\}}|\omegaD|
 \le CR^{1/D}\qquad(R>0).
\end{equation}
The same assertion holds for a finite semialgebraic chain, with absolute
multiplicities. In particular, the integral of $|\omegaD|$ is finite
on every such chain on which $f$ is bounded.
\end{theorem}
\begin{proof}
Interiority gives a $\kappa>0$ such that
$[-\kappa,\kappa]^d\subset\Newt(f)$. With
$a_* =\min_{\alpha\in\supp(f)}|a_\alpha|>0$, define
\[
 \rho_k=\max_{\alpha\in\supp(f)}|a_\alpha|2^{\pair{k}{\alpha}},
 \qquad k\in\Z^d.
\]
The support function of the contained cube gives
\begin{equation}\label{eq:rho-bound}
 \rho_k\ge a_*2^{\kappa\|k\|_1}.
\end{equation}

Partition the torus into the disjoint half-open boxes
\[
 B_k=\{x:2^{k_i}\le|x_i|<2^{k_i+1}\ (1\le i\le d)\}.
\]
On $B_k$ write $x_i=2^{k_i}z_i$, so $1\le|z_i|<2$.
Positive diagonal scaling preserves the logarithmic form. The polynomial
\[
 P_k(z)=\rho_k^{-1}z^b f(2^{k_1}z_1,\ldots,2^{k_d}z_d)
\]
has degree at most $D$ and coefficient norm exactly $1$. There is no
combination of distinct exponents in this normalization. If $|f|\le R$,
then
\begin{equation}\label{eq:scaled-sublevel}
 |P_k(z)|\le 2^{|b|_1}R/\rho_k.
\end{equation}

The rescaled sets belong to one semialgebraic family
\[
 E_\lambda=\{z:1\le|z_i|<2\text{ for every }i,
               \ (\lambda_1z_1,\ldots,\lambda_dz_d)\in S\},
 \qquad \lambda\in(0,\infty)^d.
\]
This family lies in a fixed bounded box and has dimension at most $d$.
It includes every set obtained from $S\cap B_k$. On its fixed annulus,
\[
 \omegaD=(z_1\cdots z_d)^{-1}\OmegaD,
 \qquad |\omegaD|\le|\OmegaD|.
\]
Applying Theorem~\ref{thm:sublevel} with this entire family, and then
\eqref{eq:rho-bound}, gives a constant $C_0$ independent of $k$ and $R$
with
\begin{align}
 \int_{S\cap B_k\cap\{|f|\le R\}}|\omegaD|
 &\le C_0(2^{|b|_1}R/\rho_k)^{1/D}\notag\\
 &\le C_1R^{1/D}2^{-\kappa\|k\|_1/D}.
 \label{eq:dyadic-bound}
\end{align}
Here $D=0$ would force the support of $f$ to be a singleton, contradicting
full-dimensional interiority for $d\ge1$.

Set $q=2^{-\kappa/D}\in(0,1)$. The convergent sum is explicitly
\[
 \sum_{k\in\Z^d}q^{\|k\|_1}
 =\left(1+2\sum_{j\ge1}q^j\right)^d
 =\left(\frac{1+q}{1-q}\right)^d.
\]
Countable additivity of the nonnegative integral on the disjoint boxes
proves \eqref{eq:log-integrability}. Boundary pieces are assigned to
exactly one box; they need not have zero measure in $S$. Summing the
bounds with the finitely many absolute chain multiplicities proves the
chain version.
\end{proof}

\section{Finite exceptional values and fiber periods}\label{sec:transport}
We prove the finite-exceptional-set statement by a direct scalar-gradient
argument. No Newton-polytope hypothesis is needed in this section.
Throughout, $d\ge1$ and $f\in R_d$ is nonconstant. The two points to
keep distinct are the proper geometry used for complete transport and
the later semialgebraic representatives used for integration.

\subsection{Closed embedding and the restricted differential}
Put $T=(\C^\times)^d$ and define
\begin{equation}\label{eq:closed-embedding}
 \Phi:T\longrightarrow\C^{2d},\qquad
 \Phi(z)=(z_1,\ldots,z_d,z_1^{-1},\ldots,z_d^{-1}).
\end{equation}
Its image is the closed smooth affine variety
\[
 X=\{(z,w):z_iw_i=1\ (1\le i\le d)\}.
\]
The defining differentials are independent, since their nonzero
coordinate pairs are disjoint. Replacing negative powers by powers of
$w_i$ represents $f$ by the restriction of a polynomial to $X$; write
$F:X\to\C$ for this restriction, so $F\circ\Phi=f$.

Use the Euclidean norm on $X\subset\C^{2d}$ and its induced Hermitian
metric. In torus coordinates set
\begin{equation}\label{eq:proper-radius}
 r(z)=\|\Phi(z)\|
     =\left(\sum_{i=1}^d(|z_i|^2+|z_i|^{-2})\right)^{1/2}.
\end{equation}
For every $R$, the set $\{z:r(z)\le R\}$ is compact in $T$.
Indeed it is identified with a closed bounded subset of $X$; equivalently
all $|z_i|$ are bounded above and bounded away from zero. Thus an
escaping sequence in $T$ has $r(z_n)\to\infty$ precisely when it leaves
every compact subset. The ordinary radius $\|z\|_{\C^d}$ would not have
this property.

For $x\in X$ define the \emph{restricted} differential norm
\begin{equation}\label{eq:restricted-lambda}
 \lambda(x)=\|dF_x:T_xX\longrightarrow\C\|.
\end{equation}
This is not the norm of the differential of a chosen polynomial extension
on all of $\C^{2d}$. Write $f_i=\partial f/\partial z_i$ and
$h_i(z)=1+|z_i|^{-4}$. Direct differentiation of \eqref{eq:closed-embedding}
and the weighted Cauchy--Schwarz inequality give
\begin{align}
 \|D\Phi_z(v)\|^2&=\sum_i h_i(z)|v_i|^2,
       \label{eq:embedding-metric}\\
 \lambda(\Phi(z))^2&=\sum_i\frac{|f_i(z)|^2}{h_i(z)}.
       \label{eq:lambda-formula}
\end{align}
In particular, $\lambda$ is continuous and semialgebraic.
Because $dF$ has one complex-linear output, it is either zero or
surjective as a real map to $\R^2$. Moreover its operator norm controls
both variation and the optimal right-inverse norm. Explicitly, whenever
$\lambda(\Phi(z))>0$, for every $a\in\C$ the vector
\begin{equation}\label{eq:gradient-lift}
 V_a(z)_i=\frac{a\,\overline{f_i(z)}}
                  {h_i(z)\lambda(\Phi(z))^2}
\end{equation}
satisfies
\begin{equation}\label{eq:gradient-identities}
 df_z(V_a(z))=a,\qquad
 \|D\Phi_z(V_a(z))\|=\frac{|a|}{\lambda(\Phi(z))}.
\end{equation}
These identities follow by substitution in \eqref{eq:lambda-formula};
Cauchy--Schwarz also shows that no lift of $a$ has smaller norm.
For a general map with several outputs this equality between operator
norm and least-surjectivity norm would fail. The scalar formula avoids
any need for a matrix pseudoinverse.

\subsection{Finitely many bad values}
Define
\begin{align}
 K_0&=F(\{x\in X:\lambda(x)=0\}),\label{eq:ordinary-bad}\\
 K_\infty&=\left\{c\in\C:
 \begin{array}{l}
  \exists x_n\in X,\quad \|x_n\|\to\infty,\quad F(x_n)\to c,\\
  \hfill \|x_n\|\lambda(x_n)\to0
 \end{array}\right\}.\label{eq:asymptotic-bad}
\end{align}
Here $K_0$ is the set of ordinary critical values. Every occurrence of
escape and of the differential norm in \eqref{eq:asymptotic-bad} uses
\eqref{eq:closed-embedding}--\eqref{eq:restricted-lambda}.

\begin{lemma}[Direct scalar finiteness]\label{lem:scalar-finiteness}
Both $K_0$ and $K_\infty$ are finite.
\end{lemma}
\begin{proof}
The critical locus is semialgebraic. By Standard input~\ref{std:sa},
it has a finite decomposition into connected smooth pieces, after
splitting the Nash pieces into their finitely many connected components.
The restriction of $F$ has zero differential on each such piece and is
therefore constant there, by integration along piecewise smooth paths.
Thus its image $K_0$ is finite.

For $R>1$ and $\varepsilon>0$, consider the compact semialgebraic family
\begin{equation}\label{eq:small-gradient-sphere}
 A_{R,\varepsilon}=
 \{u\in\mathbb S^{4d-1}:Ru\in X,\ R\lambda(Ru)\le\varepsilon\}.
\end{equation}
The family is compact fiberwise because $X$ is closed and $\lambda$
is continuous. It is one fixed semialgebraic family by
\eqref{eq:lambda-formula}. Lemma~\ref{lem:connecting-paths} supplies
$N\ge1$ and $L\ge1$ independent of both $R$ and $\varepsilon$.
If $C$ is a connected component and $u,v\in C$, choose a connecting
path of length at most $L$. On its arclength-parametrized arcs the
scaled path $R\gamma$ lies in $X$, so its derivative is tangent to $X$
almost everywhere and
\[
 |(F(R\gamma))'|
 \le R\lambda(R\gamma)\,\|\gamma'\|
 \le\varepsilon\|\gamma'\|.
\]
Ordinary integration on the rectifiable arcs yields
\begin{equation}\label{eq:small-gradient-diameter}
 \operatorname{diam}F(RC)\le L\varepsilon.
\end{equation}

Suppose that $K_\infty$ contains distinct $c_0,\ldots,c_N$. Put
$\eta=\min_{i\ne j}|c_i-c_j|>0$ and choose
$\varepsilon>0$ with $L\varepsilon<\eta/4$. For each $j$ let
\begin{equation}\label{eq:common-radius-set}
 E_j=\left\{R>1:
 \begin{array}{l}
 \exists x\in X,\quad\|x\|=R,\quad R\lambda(x)\le\varepsilon,\\
 \hfill |F(x)-c_j|<\eta/8
 \end{array}\right\}.
\end{equation}
It is semialgebraic, by projection, and unbounded by the definition
of $K_\infty$. Hence $E_j$ contains a whole tail $(R_j,\infty)$.
At a single radius $R>\max_jR_j$, choose the $N+1$ corresponding
points $x_j$. The normalized points $x_j/R$ lie in
$A_{R,\varepsilon}$, while
\[
 |F(x_i)-F(x_j)|>\eta-2\eta/8=3\eta/4\quad(i\ne j).
\]
Two normalized points must lie in the same one of at most $N$
components, contradicting \eqref{eq:small-gradient-diameter}.
Consequently $|K_\infty|\le N$.

This proof uses semialgebraicity of the radius sets $E_j$, not a prior
finiteness or algebraicity theorem for $K_\infty$. It does not require
curve selection at infinity or Puiseux expansions.
\end{proof}

We also write $K_0(F)=K_0$ and $K_\infty(F)=K_\infty$ when making the
function explicit.
\begin{proposition}[Finiteness of generalized critical values]\label{prop:generalized-critical}
For nonconstant $f\in R_d$, the set
\[
 B:=K_0(F)\cup K_\infty(F)\subset\C
\]
is finite.
\end{proposition}
\begin{proof}
This is the union of the two finite sets in
Lemma~\ref{lem:scalar-finiteness}.
\end{proof}

Put $U=\C\setminus B$.
\begin{lemma}[Uniform differential lower bound]\label{lem:uniform-gradient}
For every compact $Q\subset U$ there is $c_Q>0$ such that
\begin{equation}\label{eq:malgrange-uniform}
 (1+\|x\|)\lambda(x)\ge c_Q
 \qquad(x\in X,\ F(x)\in Q).
\end{equation}
\end{lemma}
\begin{proof}
An empty inverse image is harmless. Otherwise, if no positive lower
bound exists, choose $x_n$ with $F(x_n)\in Q$ and
$(1+\|x_n\|)\lambda(x_n)\to0$. Compactness gives a subsequence with
$F(x_n)\to c\in Q$. If this subsequence has a bounded subsequence,
closedness of $X$ gives a further convergent subsequence with limit
$x\in X$, $F(x)=c$, and $\lambda(x)=0$. Then $c\in K_0$.
If it has no bounded subsequence, its norms tend to infinity. The same
inequalities then give $c\in K_\infty$.
Both alternatives contradict $Q\cap B=\varnothing$.
\end{proof}

\subsection{Complete gradient transport}
We use ordinary finite-dimensional real ODE theory on an open subset of
$\C^d\simeq\R^{2d}$: a smooth vector field has unique maximal solutions,
solutions depend smoothly on initial data and finite-dimensional
parameters on their domain of existence, and a solution confined to a
compact subset of the vector-field domain cannot have a finite terminal
time. The same facts hold for time-dependent fields continuous in time and
smooth in the spatial variables, with their spatial derivatives continuous;
the relevant local Lipschitz bounds are uniform on compact time-space sets.
These are the local existence, dependence, and continuation results of
\cite[Sections 2.2 and 2.4--2.6]{Teschl}. For the local trivializations
below, dependence on finite-dimensional parameters is used only for the
explicit smooth family of autonomous fields $V_a$. The growth estimate
needed below is proved directly.

\begin{lemma}[Transport along a segment]\label{lem:complete-segment}
Let $\sigma(t)=s_0+ta$, $0\le t\le1$, be a segment contained in $U$.
For every $z_0\in f^{-1}(s_0)$, the real ODE
$\dot z=V_a(z)$, $z(0)=z_0$, is defined throughout $[0,1]$ and satisfies
$f(z(t))=\sigma(t)$. Its time-one map is a smooth diffeomorphism
$f^{-1}(s_0)\to f^{-1}(s_0+a)$. Transport of a compact cycle stays
in a compact subset of $T$ throughout this segment.
\end{lemma}
\begin{proof}
The field \eqref{eq:gradient-lift} is real smooth on
$\mathcal O=\{z\in T:\lambda(\Phi(z))>0\}$. The first identity in
\eqref{eq:gradient-identities} shows that $f(z(t))=s_0+ta$ as long as
the solution exists. Apply Lemma~\ref{lem:uniform-gradient} with the
compact set $Q=\sigma([0,1])$. Along the existing part of the solution,
\begin{equation}\label{eq:embedded-growth}
 \left\|\frac{\dd}{\dd t}\Phi(z(t))\right\|
 \le\frac{|a|}{c_Q}(1+r(z(t))).
\end{equation}
Since $r(z)\ge\sqrt{2d}>0$, its derivative exists along the smooth
solution and is at most the norm on the left. Set $C=|a|/c_Q$.
Differentiating $e^{-Ct}(1+r(z(t)))$ gives a nonpositive quantity;
hence
\begin{equation}\label{eq:gronwall-radius}
 1+r(z(t))\le(1+r(z_0))e^{Ct}\qquad(0\le t\le1)
\end{equation}
whenever the solution exists.

Choose $R_*> (1+r(z_0))e^C$. The trajectory, before any putative
terminal time at most one, lies in
\[
 \mathcal K=\{z\in T:r(z)\le R_*,\ f(z)\in Q\}.
\]
This set is compact, and \eqref{eq:malgrange-uniform} gives
$\lambda(\Phi(z))\ge c_Q/(1+R_*)>0$ on it. Thus $\mathcal K$ is a
compact subset of $\mathcal O$, not merely a bounded subset of $T$.
ODE continuation excludes such a terminal time. In particular the bound
controls both $|z_i|$ and $|z_i|^{-1}$: neither infinity nor a deleted
coordinate hyperplane can be reached in finite time.

The reversed segment uses $V_{-a}=-V_a$; uniqueness gives the inverse
of the time-one map. Smooth dependence gives smoothness of both maps.
For initial points in a compact cycle, $r(z_0)$ has a common upper
bound. The same argument is therefore uniform on that cycle, and its
sweep over $[0,1]$ is compact by continuous dependence.
\end{proof}

For $s\in\C$, write $X_s=f^{-1}(s)$.
\begin{proposition}[Complete ODE transport]\label{prop:bifurcation}
Let $B$ be the finite set in Proposition~\ref{prop:generalized-critical}
and put $U=\C\setminus B$.  Every fiber over $U$ is smooth.  Moreover,
for every piecewise $C^1$ path $\alpha:[0,1]\to U$, the minimal-norm
horizontal ODE gives a complete transport diffeomorphism
\[
 X_{\alpha(0)}\simeq X_{\alpha(1)}.
\]
In particular
\[
 f:f^{-1}(U)\longrightarrow U
\]
is locally smoothly trivial, and its fiber homology groups form a local
system on $U$.
\end{proposition}
\begin{proof}
Values outside $K_0$ have smooth fibers. First consider a $C^1$ piece
$\alpha:[u,v]\to U$; here the derivative is continuous up to the endpoints.
With $Q=\alpha([u,v])$, solve the real ODE on
$\mathcal O=\{z\in T:\lambda(\Phi(z))>0\}$ given by
\begin{equation}\label{eq:horizontal-ode}
 \dot z_i(t)=V_{\alpha'(t)}(z(t))_i
 =\frac{\alpha'(t)\,\overline{f_i(z(t))}}
        {h_i(z(t))\lambda(\Phi(z(t)))^2},
 \qquad f(z(u))=\alpha(u).
\end{equation}
The coefficients are continuous in time and smooth in space, so local
existence and uniqueness apply. By \eqref{eq:gradient-identities},
$f(z(t))=\alpha(t)$ while the solution exists. Put
$C=\sup_{[u,v]}|\alpha'|/c_Q$. Then
\begin{equation}\label{eq:linear-growth}
 \left\|\frac{\dd}{\dd t}\Phi(z(t))\right\|
 \le C(1+r(z(t))),\qquad
 1+r(z(t))\le(1+r(z(u)))e^{C(t-u)}.
\end{equation}
The second estimate follows by the same integrating-factor calculation
as \eqref{eq:gronwall-radius}. As in Lemma~\ref{lem:complete-segment},
the trajectory stays in a compact set $\{r\le R_*,\ f\in Q\}$ on which
$\lambda\ge c_Q/(1+R_*)>0$. This is a compact subset of the ODE domain;
continuation therefore covers the whole piece. The reversed path gives
the inverse transport, and smooth dependence on the initial point gives
a diffeomorphism of fibers. A piecewise $C^1$ path has a finite partition
into such pieces; composing their transports proves the path assertion.
Each lift is the minimal-norm lift by \eqref{eq:gradient-identities} and
the Cauchy--Schwarz argument following it. For a polygonal path all the
pieces use the autonomous fields of Lemma~\ref{lem:complete-segment}.

For local triviality, fix $b\in U$ and an open disk $D$ centered at $b$
with $\overline D\subset U$. For $s\in D$ let $\Theta(s,z)$ be the
time-one solution of $V_{s-b}$ starting at $z\in X_b$.
Lemma~\ref{lem:complete-segment} applies to every such segment, using the
single compact base set $\overline D$ for the differential lower bound.
Smooth dependence on $s$ and the initial point makes
\[
 \Theta:D\times X_b\longrightarrow f^{-1}(D)
\]
smooth. Its inverse sends $y$ to $f(y)$ together with the time-one
solution of $V_{b-f(y)}$ starting at $y$, and is smooth for the same
reason. The reverse-flow identity proves that the maps are inverse;
moreover $f(\Theta(s,z))=s$. Empty fibers cause no difficulty, since
reverse transport shows that a disk with empty reference fiber has
empty preimage.

The induced fiber homology groups form a local system. No canonical or
globally monodromy-invariant cycle is claimed. Uniformity on any fixed
compact cycle follows from the radius bound and continuous dependence;
no bound uniform over an entire noncompact fiber is required.
\end{proof}

\begin{remark}[Behavior at infinity is essential]\label{rem:asymptotic-example}
Newton interiority does not let one replace $B$ by $K_0$. For $c\in\C$,
consider
\[
 f_c(x,y)=c+x+\frac{(y-1)^2}{xy}.
\]
Its Newton polytope has vertices $(1,0),(-1,1),(-1,-1)$, and the origin
is their convex combination with positive coefficients $1/2,1/4,1/4$.
Along $(x,y)=(t,1)$, $t\downarrow0$, one has
\[
 f_c(t,1)=c+t,\qquad (f_c)_x(t,1)=1,\qquad (f_c)_y(t,1)=0,
\]
and, in the metric \eqref{eq:embedding-metric},
\[
 \lambda(\Phi(t,1))=\frac{t^2}{\sqrt{1+t^4}},\qquad
 r(t,1)=t+t^{-1}.
\]
Thus $r(t,1)\lambda(\Phi(t,1))\to0$ and $c\in K_\infty$.
The ordinary critical points have $y=-1$ and $x=\pm2\ii$, with critical
values $c\pm4\ii$. Hence $c$ is an ordinary regular value despite this
asymptotic degeneration.
\end{remark}
\begin{remark}[Scope and attribution of the transport proof]\label{rem:transport-background}
The normalized-gradient mechanism is a scalar specialization of the
asymptotic-critical-value approach to nonproper fibrations; see
\cite[Section 3]{Jelonek04} and \cite{KOS}. The proof here establishes the
needed finiteness by the common-radius argument and establishes
completeness by \eqref{eq:gronwall-radius}, rather than invoking a
generalized Bertini--Sard or fibration theorem. The tangency-value
alternative in \cite[Theorem 3.1]{NPP} is related background, not an
input to this proof. Appendix~\ref{app:whitney-thom} retains the independent
projective-graph/Whitney--Thom argument.

The real-smooth field $V_a$ is semialgebraic, but its flow need not be:
even $\dot u=u$ has an exponential flow. Transport provides the homology
class and local trivializations, not the semialgebraic chain used for the
endpoint estimate. Those representatives will be chosen separately in
Section~\ref{sec:primitive}.
\end{remark}

\subsection{Relative forms and holomorphic periods}
At a regular value $s$, put $X_s=f^{-1}(s)$. There is a canonical
holomorphic $(d-1)$-form $\eta_s$ on $X_s$, denoted $\omegaD/\dd f$:
locally choose a $(d-1)$-form $\eta$ with
$\dd f\wedge\eta=\omegaD$ and restrict it to the fiber. Differences
between two such choices restrict to zero. In local holomorphic
coordinates $(f,z_2',\ldots,z_d')$, this is the ordinary relative
holomorphic top form. In particular $\dd\eta_s=0$ on the fiber.

\begin{lemma}[Holomorphy of the transported period]\label{lem:holomorphy}
Let $[\gamma_s]\in H_{d-1}(X_s;\Z)$ be a locally constant homology
class over a simply connected open subset of $U$ from
Proposition~\ref{prop:bifurcation}. Then
\[
 J(s)=\int_{\gamma_s}\eta_s
\]
is well defined and holomorphic. Consequently its germ can be continued
along any path in $U$.
\end{lemma}
\begin{proof}
Closedness of $\eta_s$ and the ordinary de Rham pairing show that the
integral depends only on the homology class.  Proposition~\ref{prop:bifurcation}
provides smooth local trivializations by complete horizontal ODE transport.
Choose a compact smooth singular cycle in one fiber and transport it in a
small disc; all transported cycles remain in the image of one compact
cycle under a smooth map from a compact parameter set.

For a small oriented triangle $\Delta$ in the disc, use such a local
trivialization to sweep the cycle over $\Delta$.  This is a
$(d+1)$-chain $\Gamma_\Delta$, whose boundary is the sweep over
$\partial\Delta$.  On each edge,
$\omegaD=\dd f\wedge\eta$ and $f$ is the base coordinate, so ordinary
parametrized integration gives
\[
 \int_{\partial\Delta}J(s)\,\dd s
 =\int_{\partial\Gamma_\Delta}\omegaD
 =\int_{\Gamma_\Delta}\dd\omegaD=0.
\]
The holomorphic top form $\omegaD$ is closed.  Continuity of $J$ follows
from the same compact smooth family, and Morera's theorem proves
holomorphy.  Repeating this along overlapping discs gives continuation
along every path in $U$.  No global monodromy invariance of the transported
class is required.
\end{proof}

\section{Semialgebraic sweeps and the endpoint primitive}\label{sec:primitive}
The flows of Section~\ref{sec:transport} select the transported homology
class. They are not used as semialgebraic parametrizations. We now use
Hardt triviality again, independently of its earlier use in
Lemma~\ref{lem:connecting-paths}, to obtain the required semialgebraic sweep.
Continue to assume $0\in\Int\Newt(f)$ and $d\ge1$. Choose $b>0$
sufficiently small that $(0,b]\subset U$. Applying Hardt triviality to
the restriction of $f$ over the real interval, and decreasing $b$ once
more if necessary, gives a semialgebraic homeomorphism
\begin{equation}\label{eq:hardt}
 \Psi:(0,b]\times X_b\longrightarrow f^{-1}((0,b]),
 \qquad f(\Psi(s,z))=s,\qquad\Psi(b,z)=z.
\end{equation}
The last normalization is obtained by composing the fiber coordinate
with the inverse of its value at $b$.

Fix a compact semialgebraic cycle $\gamma_b$ representing the prescribed
transported class at $b$, and put $\gamma_s=\Psi(s,\gamma_b)$. These
classes coincide with those used for analytic continuation. Indeed,
both are locally constant in the same fiber-homology local system and
agree at $b$ on the connected interval. Thus their periods are the
function $J$ of Lemma~\ref{lem:holomorphy}.

Triangulate the finite cycle $\gamma_b$ and sweep its simplices using
\eqref{eq:hardt}. Denote the resulting finite semialgebraic $d$-chain
by $S$, and its restriction over $[a,c]\subset(0,b]$ by $S_{[a,c]}$.
Orient the base first, so
\begin{equation}\label{eq:sweep-boundary}
 \partial S_{[a,c]}=\gamma_c-\gamma_a.
\end{equation}
Shared simplex faces cancel with their cycle multiplicities. Each
truncated sweep is compact: it is the continuous image of a finite
compact cycle times $[a,c]$. Its support lies in the regular locus of
$f$. The full sweep need not have compact closure in the torus.

\begin{proposition}[Quantitative endpoint primitive]\label{prop:primitive}
There are $C>0$, $\delta>0$, and a function
$I:(0,b]\to\C$, continuous on $(0,b]$ and $C^1$ on $(0,b)$, such that
\begin{equation}\label{eq:primitive-derivative}
 I'(s)=J(s),
\end{equation}
and $I$ extends continuously to $0$ with
\begin{equation}\label{eq:primitive-endpoint}
 |I(t)-I(0)|\le Ct^\delta\quad(0<t\le b).
\end{equation}
In particular, for $0<a<c\le b$,
\begin{equation}\label{eq:sweep-integral}
 \int_a^cJ(s)\,\dd s=\int_{S_{[a,c]}}\omegaD,
\end{equation}
and, after an enlargement of $C$ if needed,
\begin{equation}\label{eq:dyadic-primitive}
 \left|\int_a^{2a}J(s)\,\dd s\right|\le C(2a)^\delta
 \quad(0<2a\le b).
\end{equation}
\end{proposition}
\begin{proof}
By Theorem~\ref{thm:log-integrability}, with $\delta=1/D$ as there,
\begin{equation}\label{eq:sweep-tail}
 \int_{|S|\cap\{|f|\le r\}}|\omegaD|\le Cr^\delta
 \qquad(r>0).
\end{equation}
This applies separately to the finitely many swept simplices and then
with their absolute multiplicities. In particular the integral over the
whole sweep is absolutely convergent, since $|f|\le b$ there. Define
\[
 I(t)=-\int_{S_{[t,b]}}\omegaD,\qquad
 I(0)=-\int_S\omegaD.
\]
The possible endpoint slices have dimension at most $d-1$ and contribute
zero to these $d$-integrals. Since $f=s$ on the sweep,
\eqref{eq:sweep-tail} proves \eqref{eq:primitive-endpoint}. It remains
to prove \eqref{eq:primitive-derivative}; the fundamental theorem of
calculus will then give \eqref{eq:sweep-integral}.

For this local Stokes calculation only, no control of escape is needed.
The following lift in the original coordinates is distinct from the
proper-metric transport field \eqref{eq:gradient-lift}. On the regular
locus of $f$, define the smooth complex vector field and smooth $(d-1)$-form
\begin{equation}\label{eq:explicit-lift}
 V=\frac{\sum_{j=1}^d\overline{\partial_j f}\,\partial/\partial x_j}
          {\sum_{j=1}^d|\partial_j f|^2},
 \qquad \eta=\iota_V\omegaD.
\end{equation}
Then $\dd f(V)=1$ and $\dd f\wedge\eta=\omegaD$; the latter follows
by contracting the identically zero $(d+1)$-form $\dd f\wedge\omegaD$.
The restriction of $\eta$ to $X_s$ is $\eta_s$.

For $0<u<v\le b$, the identity
\[
 \omegaD=\dd((f-u)\eta)-(f-u)\dd\eta
\]
and Stokes on the compact semialgebraic sweep give
\begin{equation}\label{eq:stokes-primitive}
 I(v)-I(u)=(v-u)J(v)-\int_{S_{[u,v]}}(f-u)\dd\eta.
\end{equation}
The lower boundary term is zero, and the upper one is
$(v-u)\int_{\gamma_v}\eta_v$. This is the local differentiation
identity in \cite[Lemma 9]{DvK}.

Fix $s_0\in(0,b)$. All sweeps over a sufficiently small closed
neighborhood of $s_0$ lie in one compact subset of the regular locus.
There $\dd\eta$ has bounded norm. Since $f$ takes values in $[u,v]$
on $S_{[u,v]}$, the error in the difference quotient is bounded by
\begin{equation}\label{eq:quotient-error}
 \left|\frac{I(v)-I(u)}{v-u}-J(v)\right|
 \le C_{s_0}\mass_d(S_{[u,v]}).
\end{equation}
Here mass means the sum of the Euclidean $d$-volumes of the fixed
swept pieces with their absolute multiplicities. It is finite on each
fixed compact truncation by Corollary~\ref{cor:bounded-volume}.
As $v\downarrow u=s_0$, the sets decrease to a $(d-1)$-dimensional
slice, so this mass tends to zero by continuity from above of a finite
measure. The same argument applies from the other side. Continuity of
$J$, proved independently in Lemma~\ref{lem:holomorphy}, yields
$I'(s_0)=J(s_0)$. Thus $I$ is $C^1$ on $(0,b)$. The same
compact-truncation argument at the upper endpoint gives
$I(t)\to I(b)=0$ as $t\uparrow b$.

Finally, \eqref{eq:sweep-integral} and \eqref{eq:sweep-tail}, applied
to $S_{[a,2a]}\subset\{|f|\le2a\}$, prove
\eqref{eq:dyadic-primitive}. No limiting Stokes theorem at $s=0$ has
been used: Stokes was applied only on compact truncations away from zero.
\end{proof}

\begin{corollary}[Exclusion of the forced simple pole]\label{cor:no-pole}
The period in Proposition~\ref{prop:primitive} cannot satisfy
$J(s)=A/s$ on $(0,b]$ for any constant $A\ne0$.
\end{corollary}
\begin{proof}
For $0<2a\le b$,
\[
 \int_a^{2a}\frac{A}{s}\,\dd s=A\log2.
\]
Its modulus is the positive constant $|A|\log2$, whereas
\eqref{eq:dyadic-primitive} tends to zero as $a\downarrow0$.
\end{proof}

\section{The local residue identity and minimal nonvanishing}\label{sec:residue}
By Lemmas~\ref{lem:face} and \ref{lem:chart}, the only case still
requiring proof is $d\ge1$, $0\in\Int\Newt(f)$, and
\[
 f(y)=y^{-m}u(y),\qquad m_i\ge1,\qquad u(0)\ne0.
\]
All constant terms are computed in these transformed variables, where
they equal the original constant terms. Set
\[
 \omega=\frac{\dd y_1}{y_1}\wedge\cdots\wedge\frac{\dd y_d}{y_d}.
\]

\begin{lemma}[Residue cycles near a vertex]\label{lem:residue}
For all sufficiently large $|s|$, there is a compact semialgebraic
$(d-1)$-cycle $\gamma_s\subset f^{-1}(s)$, varying locally smoothly,
such that
\begin{equation}\label{eq:residue-identity}
 R(s):=\sum_{n\ge0}\frac{\CT(f^n)}{s^{n+1}}
       =-(2\pi\ii)^{1-d}\int_{\gamma_s}\frac{\omega}{\dd f}.
\end{equation}
The series converges for all sufficiently large $|s|$.
\end{lemma}
\begin{proof}
Choose $\varepsilon>0$ so small that $u$ has no zero on the closed
polydisc $|y_i|\le\varepsilon$ and
\begin{equation}\label{eq:derivative-disc}
 |y_1\partial_1u|<m_1|u|
 \quad\text{on that polydisc}.
\end{equation}
Let $K_\varepsilon=\{|y_i|=\varepsilon\ (1\le i\le d)\}$ with its
ordered product-circle orientation. If
$|s|>\max_{K_\varepsilon}|f|$, the geometric series is uniformly
convergent on $K_\varepsilon$ and coefficient extraction gives
\begin{equation}\label{eq:cauchy-transform}
 R(s)=\frac1{(2\pi\ii)^d}
       \int_{K_\varepsilon}\frac{\omega}{s-f}.
\end{equation}
Indeed the integral of each Laurent monomial is zero unless its exponent
is zero, in which case it equals $(2\pi\ii)^d$.

For $y'=(y_2,\ldots,y_d)$ on the product circle
$K'=\{|y_j|=\varepsilon\ (j>1)\}$, solve
\begin{equation}\label{eq:roots}
 s y_1^{m_1}\prod_{j>1}y_j^{m_j}-u(y_1,y')=0.
\end{equation}
For $|s|$ uniformly large, Rouch\'e's theorem on $|y_1|=\varepsilon$
shows that there are exactly $m_1$ roots in the disc, counted with
multiplicity. None is zero because $u(0,y')\ne0$. At a root, the
$y_1$-derivative of the left side of \eqref{eq:roots} is
$(m_1u-y_1\partial_1u)/y_1$, which is nonzero by
\eqref{eq:derivative-disc}. Thus all roots are simple.

Let $\gamma_s$ be the union of these roots over $K'$. It is a compact
unbranched $m_1$-sheeted cover of $K'$, oriented by the local projections
onto the ordered base circles. Compactness follows uniformly from the
absence of roots on either $y_1=0$ or $|y_1|=\varepsilon$ and compactness
of $K'$. The defining equations and inequalities are semialgebraic.
The cover has no boundary and so defines a cycle. For $d=1$, interpret
$K'$ as a point and $\gamma_s$ as the positively weighted finite set of
roots.

As a one-variable form in $y_1$,
\[
 \frac{\dd y_1}{y_1(s-f)}
 =\frac{y_1^{m_1-1}\bigl(\prod_{j>1}y_j^{m_j}\bigr)\,\dd y_1}
 {s y_1^{m_1}\prod_{j>1}y_j^{m_j}-u(y)}.
\]
It is holomorphic at $y_1=0$. At a root $\zeta$ its residue is
$-1/(\zeta\partial_1f(\zeta,y'))$.
On this part of the fiber,
\[
 \frac{\omega}{\dd f}
 =\frac{1}{y_1\partial_1f}
   \frac{\dd y_2}{y_2}\wedge\cdots\wedge\frac{\dd y_d}{y_d}.
\]
Integrating first in $y_1$ and then over $K'$ gives
\[
 \int_{K_\varepsilon}\frac{\omega}{s-f}
 =-2\pi\ii\int_{\gamma_s}\frac{\omega}{\dd f}.
\]
Together with \eqref{eq:cauchy-transform}, this proves
\eqref{eq:residue-identity}, including its sign and normalization.
The implicit function theorem at the simple roots gives local smooth
variation of the cycle and the associated locally constant homology class.
\end{proof}

\begin{proof}[Proof of Theorem~\ref{thm:minimal}]
If the minimal-face reduction has rank zero, $f_Q$ is a nonzero constant
and its first constant term is nonzero. Otherwise apply the reduction
and unimodular change just described. Suppose for contradiction that
$\CT(f^n)=0$ for every $n\ge1$. Lemma~\ref{lem:residue} then gives
\[
 R(s)=\frac1s,\qquad
 J(s):=\int_{\gamma_s}\frac{\omega}{\dd f}
      =-\frac{(2\pi\ii)^{d-1}}s
\]
for all sufficiently large $|s|$.

Let $B$ be the finite set in Proposition~\ref{prop:bifurcation}.
Choose a large base point outside $B\cup\{0\}$ and continue the
homology class of the residue cycle along a path in
$\C\setminus(B\cup\{0\})$ to a sufficiently small positive $b$.
Such a path exists because the complement of finitely many points in
$\C$ is path connected. Lemma~\ref{lem:holomorphy} continues the scalar
period. Since the rational function $-(2\pi\ii)^{d-1}/s$ is holomorphic
along the same path, uniqueness of analytic continuation gives
\[
 J(s)=-\frac{(2\pi\ii)^{d-1}}s\quad(0<s\le b).
\]
Here $J$ denotes the period of the continued homology class, not the
integral over the original fixed product torus after leaving its initial
convergence region. The cycle may have monodromy; the displayed scalar
identity is what is continued.

Proposition~\ref{prop:primitive} applies to this class and its
semialgebraic representatives on $(0,b]$. Corollary~\ref{cor:no-pole}
contradicts the displayed identity, since $(2\pi\ii)^{d-1}\ne0$.
The assumed universal vanishing is therefore impossible.
\end{proof}

\section{Infinite nonvanishing and the torus Mathieu property}\label{sec:corollaries}
\begin{proof}[Proof of Theorem~\ref{thm:main}]
Fix $M\ge0$ and set $q=M+1$. By Lemma~\ref{lem:powers},
$0\in\Newt(f^q)$. Theorem~\ref{thm:minimal} gives $k\ge1$ with
\[
 0\ne\CT((f^q)^k)=\CT(f^{qk}).
\]
Since $qk\ge q>M$, this proves the assertion. In particular, no separate
argument concerning higher-order poles is required.
\end{proof}

\begin{corollary}[Newton-polytope classification]\label{cor:classification}
For $0\ne f\in R_d$,
\[
 \CT(f^n)=0\text{ for every }n\ge1
 \quad\Longleftrightarrow\quad
 0\notin\Newt(f).
\]
\end{corollary}
\begin{proof}
One direction is the contrapositive of Theorem~\ref{thm:minimal}.
For the other, strict separation of a point and a compact convex set
gives a linear functional $\ell$ and $\delta>0$ with
$\ell(\alpha)\ge\delta$ for every $\alpha\in\supp(f)$. Every exponent
in $f^n$ is a sum of $n$ such exponents, and has $\ell$-value at least
$n\delta>0$. The zero exponent cannot occur.
\end{proof}

\begin{corollary}[Laurent-polynomial Mathieu property]\label{cor:mathieu}
Let $f,h\in R_d$. If $\CT(f^n)=0$ for every $n\ge1$, then
$\CT(hf^n)=0$ for all sufficiently large $n$.
\end{corollary}
\begin{proof}
The assertion is immediate if $f=0$ or $h=0$. Otherwise use
Corollary~\ref{cor:classification} and choose $\ell,\delta$ as in its
proof. Put $A_h=\min_{\beta\in\supp(h)}\ell(\beta)$. Any exponent in
$hf^n$ has $\ell$-value at least $A_h+n\delta$, which is positive for
all sufficiently large $n$. Thus its constant coefficient is zero.
\end{proof}

\begin{corollary}[Compact tori]\label{cor:torus}
Let $K=\T^d$ with normalized Haar measure. For finite Fourier sums
$F,H$ on $K$, if $\int_K F^n=0$ for every $n\ge1$, then
$\int_K HF^n=0$ for all sufficiently large $n$.
\end{corollary}
\begin{proof}
Identify a character with $x^\alpha=e^{\ii\pair{\alpha}{\theta}}$, $\alpha\in\Z^d$.
Direct integration in each circle coordinate gives
\[
 \frac1{(2\pi)^d}\int_{[0,2\pi]^d}
 e^{\ii\pair{\alpha}{\theta}}\,\dd\theta
 =\begin{cases}1,&\alpha=0,\\0,&\alpha\ne0.\end{cases}
\]
Hence the Haar integral of any Laurent polynomial restricted to $K$
is its constant coefficient. Apply Corollary~\ref{cor:mathieu}.
\end{proof}

\begin{remark}
Finite Fourier sums are exactly the continuous $K$-finite functions on
a torus. One direction follows because each character spans a
one-dimensional translation module. Conversely, on a finite-dimensional
translation-invariant space, average an inner product over $K$ and
simultaneously diagonalize the commuting unitary translations. An
eigenfunction is a scalar multiple of the corresponding character,
by evaluation at the identity. The continuous characters of $\T^d$
are precisely those indexed by $\Z^d$. This is the usual torus
interpretation of \cite{DvK}; no assertion about nonabelian groups is
needed in this paper.
\end{remark}

\section{Provenance and mathematical scope}\label{sec:scope}
The theorem proved is the arbitrary-rank nonvanishing theorem of
Duistermaat and van der Kallen. The residue construction in
Lemma~\ref{lem:residue} is their local construction, and the calculation
\eqref{eq:stokes-primitive} is their primitive-differentiation identity.
The contribution proposed for independent examination is the uniform
middle-dimensional sublevel estimate and its dyadic use in
Theorem~\ref{thm:log-integrability}. This replaces the argument making
the Haar form regular over finite fibers of a smooth compactification.

The semialgebraic integration estimates in \cite{HKT} are antecedents,
not conclusions attributed to the present paper. The coefficient-derivative
induction in Section~\ref{sec:sublevel} is also naturally compared with
geometric sublevel-set estimates; see, for example, \cite{Gressman} and
the references there. The specific estimate here concerns the modulus
of a restricted holomorphic top form and is uniform in a fixed
semialgebraic family. These precise hypotheses, rather than merely
similar terminology, must be used in any priority comparison.

Question~9.1 of \cite{Kiselman} motivates the search for a simpler
nonvanishing proof. The main route retains substantial, explicitly
identified semialgebraic and integration results, together with ordinary
finite-dimensional ODE theory. Whitney stratification and Thom isotopy
are confined to the independent appendix proof and are not dependencies
of the main route. The argument is therefore described as resolution-free
rather than as elementary in every respect. The stronger singularity and growth conclusions of
\cite[Theorem 5]{DvK} are outside the claims made here. A proof and
generalization using Newton--Puiseux methods in one variable is discussed
in \cite{ES}; the present theorem has no rank restriction.

The intended future formalization has one core endpoint,
Theorem~\ref{thm:minimal}. Lemma~\ref{lem:powers} and the separation
argument then give Theorem~\ref{thm:main} and
Corollary~\ref{cor:mathieu}. The semialgebraic and integration inputs in
Section~\ref{sec:sa} and the ODE facts identified in
Section~\ref{sec:transport} remain genuine mathematical dependencies:
a formal proof must establish them or invoke proved library theorems,
not add them as unproved project axioms. In particular, the connecting-path
lemma is derived from the existing uniform-volume bound; finiteness of
$K_\infty$ is then proved rather than supplied by an external theorem.
Formalization of the independent appendix proof would additionally need
Standard input~\ref{std:whitney}; that input is not required to formalize
the main scalar-gradient route.

\subsection*{Acknowledgment of computational assistance}
ChatGPT assisted with the development and drafting of this manuscript.
Comments attributed to Claude informed revisions. These tools and their
informal checks do not constitute independent mathematical verification
or a Lean formalization.

\appendix
\section{Whitney--Thom alternative for finite bifurcation and cycle transport}\label{app:whitney-thom}
This appendix records an independent proof of the finite-exceptional-set
and topological transport conclusion needed for the period argument. It
is not used in the main proof. The exceptional set below need not equal
$K_0\cup K_\infty$. This route does not assert the quantitative gradient
bounds or the specific transport diffeomorphisms of
Proposition~\ref{prop:bifurcation}.

\begin{standard}[Complex algebraic stratification and first isotopy]\label{std:whitney}
A complex algebraic variety, together with finitely many specified
algebraic subsets, admits a finite complex algebraic Whitney
stratification compatible with those subsets. If a map from a closed
Whitney-stratified subset of a smooth manifold to a smooth base is proper
and restricts to a submersion on every stratum, it is locally topologically
trivial by homeomorphisms preserving the strata.
\end{standard}

These are the algebraic Whitney-stratification theorem and Thom's first
isotopy lemma; see \cite{Kaloshin,Mather}. For a direct treatment of
nonproper polynomial maps, see \cite{DJ}, especially Theorems 3.1 and
5.1 of its arXiv version.

\begin{proposition}[Whitney--Thom finite exceptional set]\label{prop:bifurcation-whitney}
For a nonconstant $f\in R_d$, there is a finite set $B_{\mathrm W}\subset\C$
such that all fibers over $U_{\mathrm W}=\C\setminus B_{\mathrm W}$ are
smooth and
\[
 f:f^{-1}(U_{\mathrm W})\longrightarrow U_{\mathrm W}
\]
is locally topologically trivial.
\end{proposition}
\begin{proof}
Let $\overline\Gamma$ be the closure of the graph of $f$ in
$(\Pj^1)^d\times\Pj^1$, let $p:\overline\Gamma\to\Pj^1$ be projection,
and let $D_\infty$ be the complement of the open graph identified with
$\torus$. The projection $p$ is proper. Choose a finite complex
algebraic Whitney stratification compatible with $D_\infty$.

For a stratum $W$, the critical locus of $p|_W$ is algebraic in $W$.
At such a point its complex differential has rank zero, because the
base has complex dimension one. On the smooth locus of any irreducible
component of this critical locus, the restriction of $p$ consequently
has zero differential and is locally constant. It is therefore constant
on that component. Equivalently, a rational function with zero
differential on an irreducible complex variety is constant. There are
only finitely many components and strata, so their critical images form
a finite subset of $\Pj^1$.

Remove that finite subset and the value $\infty$. On the remaining base,
$p$ is a proper stratified submersion, so Standard
input~\ref{std:whitney} gives local trivializations preserving all strata.
They preserve $D_\infty$ and its complement. Restricting to the latter
gives the claimed local trivializations of $f$. At a point of this open
graph, submersivity on its stratum implies submersivity on the ambient
smooth torus; hence every remaining fiber is smooth. Empty fibers, if
present, cause no difficulty in this statement. In the application a
nonempty fiber is supplied by the explicit residue construction.
\end{proof}

\begin{remark}
The adjective \emph{complex algebraic} matters in this appendix. Real
semialgebraic strata can map to curves of real critical values. Such a
real exceptional set need not leave a connected complement through which
the required analytic germ can be transported. The projective graph may
be singular; no resolution theorem is hidden in this argument.
\end{remark}

\begin{remark}
Proposition~\ref{prop:bifurcation-whitney} can replace
Propositions~\ref{prop:generalized-critical}--\ref{prop:bifurcation} in
the nonvanishing proof.  For the period itself, one then uses the local
compact-cycle argument: on a sufficiently small disc
of regular values, lift the two real coordinate vector fields only near the
compact cycle support, transport the finite cycle smoothly, and apply the
triangle-sweep/Morera argument of Lemma~\ref{lem:holomorphy}.  Thus the
appendix route supplies the same analytic continuation without any appeal to
the ODE completeness argument of Section~\ref{sec:transport}.
\end{remark}

\begin{thebibliography}{99}
\bibitem{BCR}
J.~Bochnak, M.~Coste, and M.-F.~Roy,
\emph{Real Algebraic Geometry}, Ergebnisse der Mathematik und ihrer
Grenzgebiete (3), vol.~36, Springer, Berlin, 1998.

\bibitem{DJ}
S.~T.~Dinh and Z.~Jelonek,
Thom isotopy theorem for nonproper maps and computation of sets of
stratified generalized critical values,
\emph{Discrete Comput. Geom.} \textbf{65} (2021), 279--304.
\href{https://doi.org/10.1007/s00454-019-00087-w}{doi:10.1007/s00454-019-00087-w}.
See also \arxiv{1805.02681}, version~2 (2018), Theorems~3.1 and~5.1.

\bibitem{DvK}
J.~J.~Duistermaat and W.~van der Kallen,
Constant terms in powers of a Laurent polynomial,
\emph{Indag. Math. (N.S.)} \textbf{9} (1998), no.~2, 221--231.
\href{https://doi.org/10.1016/S0019-3577(98)80020-7}{doi:10.1016/S0019-3577(98)80020-7}.

\bibitem{ES}
A.~van den Essen and J.~Schoone,
Kernels of linear maps: a generalization of Duistermaat and van der
Kallen's theorem, \arxiv{2305.10062} (2023).

\bibitem{Gressman}
P.~T.~Gressman,
Uniform geometric estimates for sublevel sets,
\arxiv{0909.0875} (2009).

\bibitem{HKT}
M.~Hanamura, K.~Kimura, and T.~Terasoma,
Integrals of logarithmic forms on semi-algebraic sets and a generalized
Cauchy formula, Part I: convergence theorems,
\arxiv{1509.06950}, version~1 (2015).
In particular, Sections~1--2, Propositions~2.5 and~2.8,
and Theorem~2.9.

\bibitem{Hardt}
R.~M.~Hardt,
Semi-algebraic local-triviality in semi-algebraic mappings,
\emph{Amer. J. Math.} \textbf{102} (1980), no.~2, 291--302.
\href{https://doi.org/10.2307/2374240}{doi:10.2307/2374240}.

\bibitem{Kaloshin}
V.~Kaloshin,
A geometric proof of the existence of Whitney stratifications,
\arxiv{math/0010144} (2000), Theorem~1.

\bibitem{Kiselman}
C.~O.~Kiselman,
Questions inspired by Mikael Passare's mathematics,
\emph{Afr. Mat.} \textbf{25} (2014), 271--288;
published online October~24, 2012.
\href{https://doi.org/10.1007/s13370-012-0107-5}{doi:10.1007/s13370-012-0107-5}.
Section~9, Question~9.1.

\bibitem{Mather}
J.~Mather,
Notes on topological stability,
\emph{Bull. Amer. Math. Soc. (N.S.)} \textbf{49} (2012), no.~4, 475--506.

\bibitem{OS}
T.~Ohmoto and M.~Shiota,
$C^1$-triangulations of semialgebraic sets,
\emph{J. Topol.} \textbf{10} (2017), no.~3, 765--775.
\href{https://doi.org/10.1112/topo.12024}{doi:10.1112/topo.12024}.
See also \arxiv{1505.03970}, version~2 (2017), especially Section~4.

\bibitem{Jelonek04}
Z.~Jelonek,
On asymptotic critical values and the Rabier theorem,
\emph{Banach Center Publ.} \textbf{65} (2004), 125--133.
\href{https://doi.org/10.4064/bc65-0-9}{doi:10.4064/bc65-0-9}.

\bibitem{KOS}
K.~Kurdyka, P.~Orro, and S.~Simon,
Semialgebraic Sard theorem for generalized critical values,
\emph{J. Differential Geom.} \textbf{56} (2000), no.~1, 67--92.
\href{https://doi.org/10.4310/jdg/1090347525}{doi:10.4310/jdg/1090347525}.
In particular, Theorem~3.3 and Section~3.4.

\bibitem{NPP}
T.~T.~Nguyen, P.-P.~Pham, and T.-S.~Pham,
Bifurcation sets and global monodromies of Newton non-degenerate
polynomials on algebraic sets, \arxiv{1803.09654} (2018).
See Theorem~3.1 for the tangency-value approach.

\bibitem{Teschl}
G.~Teschl,
\emph{Ordinary Differential Equations and Dynamical Systems},
Graduate Studies in Mathematics, vol.~140, American Mathematical Society,
Providence, RI, 2012. Sections~2.2 and~2.4--2.6.

\end{thebibliography}
\end{document}